% !TEX root = ./main.tex

\ustcsetup{
  title                = {市场心智：基于 LLM 驱动的\\多源信息聚合与市场共识早期发现\\及其金融实践},
  title*               = {MarketMind: LLM-based Multi-Source Integration and Market Consensus Discovery with Financial Practice},
  author               = {彭晗},
  author*              = {Han Peng},
  speciality           = {计算机科学与技术},
  speciality*          = {Computer Science and Technology},
  supervisor           = {邢凯~教授},
  supervisor*          = {Prof. Kai Xing},
  % practice-supervisor  = {XXX~教授, XXX~教授},  % 专业/工程学位的实践导师
  % practice-supervisor* = {Prof. XXX, Prof. XXX},
  % date                 = {2025-01-01},  % 完成时间，默认为今日
  %
  % department           = {数学科学学院},  % 院系，本科生需要填写
  student-id           = {PB22111636},  % 学号，本科生需要填写
  %
  % secret-level         = {秘密},     % 绝密|机密|秘密|控阅，注释本行则公开
  % secret-level*        = {Secret},  % Top secret | Highly secret | Secret
  % secret-year          = {10},      % 保密/控阅期限
  %
  % 数学字体
  % math-style           = GB,  % 可选：GB, TeX, ISO
  math-font            = xits,  % 可选：stix, xits, libertinus
}


% 加载宏包

% 定理类环境宏包
\usepackage{amsthm}

% 插图
\usepackage{graphicx}

% 英文图题、表题
\usepackage{bicaption}

% 三线表
\usepackage{booktabs}

% 表注
\usepackage{threeparttable}

% 跨页表格
\usepackage{longtable}

% 算法
\usepackage[ruled,linesnumbered]{algorithm2e}

% SI 量和单位
\usepackage{siunitx}

% 参考文献使用 BibTeX + natbib 宏包
% 顺序编码制
\usepackage[sort]{natbib}
\bibliographystyle{ustcthesis-numerical}

% 著者-出版年制
% \usepackage{natbib}
% \bibliographystyle{ustcthesis-authoryear}

% 本科生参考文献的著录格式
% \usepackage[sort]{natbib}
% \bibliographystyle{ustcthesis-bachelor}

% 参考文献使用 BibLaTeX 宏包
% \usepackage[style=ustcthesis-numeric]{biblatex}
% \usepackage[bibstyle=ustcthesis-numeric,citestyle=ustcthesis-inline]{biblatex}
% \usepackage[style=ustcthesis-authoryear]{biblatex}
% \usepackage[style=ustcthesis-bachelor]{biblatex}
% 声明 BibLaTeX 的数据库
% \addbibresource{bib/ustc.bib}

% 配置图片的默认目录
\graphicspath{{figures/}}

% 数学命令
\makeatletter
\newcommand\dif{%  % 微分符号
  \mathop{}\!%
  \ifustc@math@style@TeX
    d%
  \else
    \mathrm{d}%
  \fi
}
\makeatother
\newcommand\eu{{\symup{e}}}
\newcommand\iu{{\symup{i}}}

% 用于写文档的命令
\DeclareRobustCommand\cs[1]{\texttt{\char`\\#1}}
\DeclareRobustCommand\env[1]{\texttt{#1}}
\DeclareRobustCommand\pkg[1]{\textsf{#1}}
\DeclareRobustCommand\file[1]{\nolinkurl{#1}}

% hyperref 宏包在最后调用
\usepackage{hyperref}

\IfFontExistsTF{Times New Roman}{%
  \newfontfamily\ustctimes{Times New Roman}%
}{%
  \newfontfamily\ustctimes{TeX Gyre Termes}%
}

\renewcommand\contentsname{\texorpdfstring{目\quad 录}{目录}}
\renewcommand\bibname{\texorpdfstring{参\quad 考\quad 文\quad 献}{参考文献}}

\makeatletter

% 正文：小四宋体，英文和数字使用 Times New Roman，固定行距 20 磅，首行缩进 2 字符。
\AtBeginDocument{%
  \fontsize{12bp}{20bp}\selectfont
  \setlength{\parskip}{0bp}%
  \ctexset{autoindent=2}%
}

% 章、节、段标题与目录采用本科论文格式规范中的中文编号体系。
\ctexset{
  chapter = {
    name       = {第,章},
    number     = \chinese{chapter},
    format     = \centering\heiti\fontsize{16bp}{20bp}\selectfont,
    aftername  = \quad,
    beforeskip = 24bp,
    afterskip  = 18bp,
  },
  section = {
    name       = {第,节},
    number     = \chinese{section},
    format     = \heiti\fontsize{15bp}{20bp}\selectfont,
    aftername  = \quad,
    beforeskip = 24bp,
    afterskip  = 6bp,
  },
  subsection = {
    number     = \chinese{subsection},
    format     = \heiti\fontsize{14bp}{20bp}\selectfont,
    aftername  = {、},
    indent     = 0bp,
    beforeskip = 12bp,
    afterskip  = 6bp,
  },
}

% 摘要：标题三号黑体；内容小四宋体/Times New Roman，固定行距 20 磅。
\renewenvironment{abstract}{%
  \ustcsetup{language=chinese}%
  \clearpage
  \phantomsection
  \pdfbookmark[0]{摘要}{abstract.zh}%
  {\centering\heiti\fontsize{16bp}{16bp}\selectfont 摘\quad 要\par}%
  \vspace{18bp}%
  \fontsize{12bp}{20bp}\selectfont
}{%
  \par\vspace{20bp}%
  \begingroup
    \setlength{\hangindent}{4em}%
    \noindent\heiti 关键词：\songti\ustc@keywords@text\par
  \endgroup
  \ustc@reset@main@language
}

\renewenvironment{abstract*}{%
  \ustcsetup{language=english}%
  \clearpage
  \phantomsection
  \pdfbookmark[0]{ABSTRACT}{abstract.en}%
  {\centering\ustctimes\bfseries\fontsize{16bp}{16bp}\selectfont ABSTRACT\par}%
  \vspace{18bp}%
  \ustctimes\fontsize{12bp}{20bp}\selectfont
}{%
  \par\vspace{20bp}%
  \begingroup
    \setlength{\hangindent}{6.5em}%
    \noindent\ustctimes\bfseries Key Words:\space
    \normalfont\ustctimes\ustc@keywords@en@text\par
  \endgroup
  \ustc@reset@main@language
}

% 目录：小四宋体，固定行距 20 磅；章左顶格加粗，节缩进 2 字符，段缩进 4 字符。
\titlecontents{chapter}
  [0bp]{\fontsize{12bp}{20bp}\selectfont\heiti}
  {\contentspush{\thecontentslabel\quad}}{}
  {\songti\ustc@leaders\thecontentspage}
\titlecontents{section}
  [2em]{\fontsize{12bp}{20bp}\selectfont}
  {\contentspush{\thecontentslabel\quad}}{}
  {\ustc@leaders\thecontentspage}
\titlecontents{subsection}
  [4em]{\fontsize{12bp}{20bp}\selectfont}
  {\contentspush{\thecontentslabel、}}{}
  {\ustc@leaders\thecontentspage}

\makeatother
